\section{Implikationen für UX-Strategien in Organisationen}
\label{sec:06-03_organizational-ux-implications}

Die empirischen Befunde dieser Arbeit zeigen, dass \Gls{uxstrategien} im organisationalen Kontext nicht isoliert als methodisches Vorgehen, sondern als integrierter Gestaltungs- und Steuerungsrahmen verstanden werden müssen. Ausgehend von den Ergebnissen in \hyperref[ch:05_analysis-and-results]{Kapitel~5} und deren theoretischer Einordnung in \hyperref[sec:06-01_theoretical-positioning]{Abschnitt~6.1} lassen sich daraus mehrere praxisrelevante Implikationen für die strategische Gestaltung interner digitaler Plattformen ableiten.

Ein zentraler Befund ist die Bedeutung einer klaren Struktur. Die Effizienz- und Orientierungsgewinne durch eine konsistente \Gls{ia} zeigen, dass eine gut aufgebaute Informationsstruktur wesentlich zum langfristigen Nutzen organisationaler Systeme beiträgt. \Gls{ia} sollte daher als kontinuierliche Gestaltungsaufgabe verstanden und regelmäßig weiterentwickelt werden. Strategische \Gls{ux}-Roadmaps können hierzu Ziele wie die Reduktion kognitiver Belastung, einheitliche Navigationslogiken und klar erwartungskonforme Kategorien festlegen \cite{303:Rosenfeld2015,108:ArangoFirstPrinciples2024,113:NNGIntranetGuidelines}. Regelmäßige \Gls{usability}-Tests und Nutzungsanalysen helfen dabei, strukturelle Schwächen frühzeitig zu erkennen und Qualitätsverluste zu vermeiden.

Die hohe Wirksamkeit von Einstiegs- und Übersichtsseiten zeigt ihr strategisches Potenzial als aktive Steuerungselemente des Nutzungserlebnisses. Statt sie nur als Informationssammlung zu betrachten, sollten sie gezielt an zentralen Nutzungsszenarien und organisationalen Kernprozessen ausgerichtet und regelmäßig an veränderte Arbeitsrealitäten angepasst werden. Dadurch lassen sich Onboarding-Aufwände reduzieren und Effizienzgewinne für erfahrene Nutzende erzielen \cite{307:DWGFindability,113:NNGIntranetGuidelines}. Eine klare Priorisierung relevanter Inhalte unterstützt zudem die Entscheidungsfindung und senkt Suchaufwände im Arbeitsalltag.

Im Hinblick auf Rollenpassung und Nutzerorientierung zeigen die Ergebnisse, dass hohe Gebrauchstauglichkeit nicht zwingend durch individualisierte Oberflächen erreicht werden muss. Eine gute Passung zu unterschiedlichen Rollen entsteht vielmehr durch vollständige Inhalte, transparente Zuständigkeiten und konsistente Strukturen. \Gls{personas} entfalten ihren strategischen Wert daher weniger als Grundlage technischer Personalisierung, sondern als gemeinsames Referenzmodell für Design-, Struktur- und Inhaltsentscheidungen \cite{112:StepTwoPersonas,308:DWGSegmentation}. Auf diese Weise unterstützen sie die Priorisierung von Inhalten und eine kohärente Gesamtstruktur, ohne die \Gls{ia} durch zu starke Segmentierung zu zersplittern.

Da \Gls{uxstrategien} vor allem über wahrnehmbare Gestaltung und Nutzungserfahrungen wirken, sollte \Gls{ux} stärker als messbare Ergebnisqualität verstanden werden. Kennzahlen wie Auffindbarkeit, Bearbeitungszeit oder Fehlerraten ermöglichen eine kontinuierliche Bewertung und datenbasierte Entscheidungen \cite{300:FraunhoferIESE2018,101:Khan2021}. So wird auch der Beitrag von \Gls{ux}-Investitionen zur Produktivität und Stabilität von Arbeitsprozessen sichtbar \cite{301:Infragistics2015,100:WhartonUX2015}.

Gleichzeitig machen die identifizierten Akzeptanzbarrieren deutlich, dass \Gls{uxstrategien} eng mit Governance-, Verantwortungs- und Kommunikationsstrukturen verknüpft sind. Klare Pflegeprozesse, verbindliche Qualitätsstandards und transparente Rollenmodelle sind Voraussetzungen dafür, dass gute \Gls{usability} langfristig wirksam bleibt \cite{115:NNGContentManagement,309:NNGContentGovernance,116:StepTwoGuidance}.

Barrierefreiheit und Mehrsprachigkeit sind als eigenständige Qualitätsdimensionen organisationaler \Gls{uxstrategien} zu berücksichtigen, da sie den gleichberechtigten Zugang zu Informationen und die langfristige Nutzbarkeit digitaler Plattformen wesentlich beeinflussen. Anforderungen an Barrierefreiheit lassen sich über ergonomische Gestaltungsleitlinien für zugängliche Software der \acrshort{iso}~9241-171 sowie über internationale Richtlinien für barrierefreie Webinhalte (\Gls{wcag}) und den europäischen Referenzstandard EN~301~549 konkretisieren \cite{502:ISO9241-171,503:WCAG22,504:EN301549}. Entsprechend sollten diese Anforderungen in Designrichtlinien, Content-Standards und technischen Vorgaben verbindlich verankert und über klar definierte Verantwortlichkeiten in die Governance integriert werden.
Mehrsprachigkeit ist nicht auf reine Übersetzungsprozesse zu reduzieren, sondern erfordert konsistente Struktur- und Benennungsmodelle, etwa durch systematische Terminologiepflege sowie eine eindeutige Zuordnung von Labels und Metadaten je Sprache. Empfehlungen des \Gls{w3c} betonen zudem die Bedeutung korrekter Sprachdeklaration und sprachbezogener Auszeichnung für Verarbeitung, Suche und Darstellung in mehrsprachigen Systemen \cite{505:W3Ci18nAuthoring,506:W3Ci18nQuickTips}. Eine vertiefte empirische Analyse von Barrierefreiheit und Mehrsprachigkeit war nicht Gegenstand der vorliegenden Arbeit, stellt jedoch einen relevanten Ansatzpunkt für weiterführende Untersuchungen dar.

Ergänzend sind gezielte Einführungs- und Kommunikationsmaßnahmen notwendig, um neue Systeme in bestehende Arbeitsabläufe zu integrieren und klare Nutzungserwartungen zu schaffen \cite{110:Lewin1947,418:Kotter2012,118:NNGUserInvolvement}. \Gls{ux}-Strategie wird damit zu einer organisationsweiten Gestaltungsaufgabe, die technische, inhaltliche und kulturelle Aspekte miteinander verbindet.

Schließlich unterstreichen die Ergebnisse die Bedeutung kontinuierlicher Evaluation. Die Kombination aus quantitativen \Gls{usability}-Metriken und qualitativen Rückmeldungen ermöglicht eine evidenzbasierte Weiterentwicklung der Plattform und unterstützt eine lernende Organisationskultur \cite{121:NNGUsabilityTesting,124:Triangulation2023}. So können strukturelle Defizite frühzeitig erkannt und Anpassungen systematisch umgesetzt werden.

Insgesamt zeigen die Implikationen, dass erfolgreiche \Gls{uxstrategien} weniger durch einzelne Designmaßnahmen als durch die konsequente Verzahnung von Strukturqualität, Governance, Kommunikation und kontinuierlicher Evaluation getragen werden. \Gls{ux} wird damit zu einem strategischen Treiber organisationaler Leistungsfähigkeit und Wissensarbeit, der langfristig geplant, gesteuert und institutionell verankert werden muss \cite{417:Benyon2019,403:Westerman2014}.
